\chapter{Introduction}

Particle physics experiments generate data at unprecedented scales, with the Large Hadron Collider (LHC) producing approximately one petabyte of raw collision data per second~\cite{cern_computing}. To extract meaningful physics from this stream, a two-staged trigger system is applied. The first level (L1), consisting of custom hardware processes, provides coarse-grain filtering, selecting events at a rate of $100kHz$, and filtering-out most of them. The second level (L2), running custom software for event reconstruction, samples events at just $400Hz$ for offline storage and processing~\cite{cms_trigger_2017,cms_hlt_run2}. This work focuses on second-level filtering and addresses a particularly challenging reconstruction problem: identifying and characterizing $\tau$ leptons decaying to three charged pions ($\tau \to \pi\pi\pi \nu_\tau$)~\cite{pdg_tau_branching} in the low transverse momentum ($p_T$) regime, where conventional algorithms exhibit degraded performance~\cite{cms_tau_performance_2018,cms_deeptau_v25,abcnet-low-pt}.

\section{Motivation}

The problem of jet tagging for various particles has seen multiple advancements over the last years, with models such as ParticleNet~\cite{particlenet}, ParticleTransformer~\cite{particle_transformer_2022}, and OmniLearned~\cite{omnilearned}. However, in their base implementation, these models primarily focus on jets with medium-to-high transverse momentum ($p_T > 500$~GeV). On the other hand, in the low-$p_T$ regime, the fundamental assumption of jet-based seeding - approach when we start the process of $\tau$ identification from finding a jet of its decay products - breaks down.

\subsubsection{Failure of Jet-Seeded Reconstruction}

As demonstrated in~\cite{abcnet-low-pt}, the maximum $\Delta R = \sqrt{\left(\Delta\eta\right)^2 + \left(\Delta\varphi\right)^2}$ spread between any pair of the three charged pions in $\tau \to \pi\pi\pi\nu$ decays exceeds 0.4 (the standard jet cone size) for visible $\tau_h$ with $p_T$ below approximately 8~GeV and can reach values above 1.0 at $p_T < 5$~GeV~\cite{atlas_low_pt_tracking_2016}. This geometric spread arises from the non-negligible $\tau$ mass: at low momentum, decay products are no longer Lorentz-boosted into a collimated cone but instead diverge across a large angular region. Consequently, the probability of containing all three charged pions within a single ak4 jet radius drops to nearly 0\% below 3~GeV, rendering typical reconstruction algorithms ineffective.

\subsubsection{The CMS Low-$p_T$ $\tau_h$ Algorithm}

To address physics analyses requiring $\tau_h$ identification below 10~GeV, such as B-physics measurements ($B_c \to J/\psi\tau\nu$)~\cite{lhcb_rd_semitauonic_2019,belle_b0_kstar_tautau} and low mass dark matter searches CMS developed a dedicated seedless algorithm targeting the 3-prong decay channel~\cite{abcnet-low-pt}. The reconstruction proceeds in two stages: selecting charged pion candidates near the preselected proton collision point ($p_T > 0.5$~GeV, $|\eta| < 2.5$); followed by ABCNet~\cite{abcnet} classification of $\tau_h$ candidates preselected based on the quality of reconstructed $\tau$ candidate origin and decay vertex. However, on average, the algorithm manages to recover the entire triplet of pions for less than 40\% of events, even for $\tau$ leptons with $p_T$ approaching 20~GeV~ - the higher end of the low-$p_T$ regime.\\

This work aims to replace the existing low-$p_T$ $\tau_h$ reconstruction algorithm at CMS, whose limited recovery rate falls well short of what downstream physics analyses require. Architectures that can identify the correct triplet and reconstruct the $\tau$-lepton from small candidate sets ($<50$ elements) already exist~\cite{tau_unified_ml_2024,gnn_tau_future_colliders}; reaching such a set, however, requires first refining a much larger initial pool ($>600$ pion candidates). It is precisely this refinement step in the low-$p_T$ regime that the literature has left underexplored, and that this thesis takes as its central problem — addressing simultaneously a concrete shortcoming of CMS reconstruction and a methodological gap in candidate-refinement for sparse, low-$p_T$ signatures.

\section{Contributions}

We prepared a dedicated low-$p_T$ $\tau_h$ dataset of neutral B-meson decays into $\tau$-pairs, generated by the standard CMS Monte Carlo chain (Pythia~8~\cite{10.21468/SciPostPhysCodeb.8} + EvtGen~\cite{EvtGen} + CMSSW~14~\cite{CMSSW}) and converted into a training-ready Parquet format with detector-level features and per-track provenance labels. The underlying ROOT samples remain private to CERN; this thesis releases only the conversion pipeline, derived statistics, and aggregate descriptions of the data.

We developed a two-stage cascade reconstruction model directly inspired by the ABCNet low-$p_T$ pipeline, but replacing its analytic preselection — the $\tau_h$ candidate construction and quality cuts that already discard most of the recoverable signal — with a dedicated learned pre-filter. The pre-filter ranks all reconstructed pion candidates in an event ($>1100$ on average) and forwards a refined shortlist to a Particle-Transformer-based reranker, which operates only on detector-level quantities at inference time. The two stages are complementary by design: Stage 1 is a per-track scorer, cheap enough to apply to the full event but unable to exploit inter-track structure on its own, while Stage 2 is a self-attention reranker that captures the missing pairwise structure but whose $O(N^2)$ cost makes it prohibitively expensive on the full $>1100$-track input.

To make the Stage-2 attention sensitive to $\tau \to 3\pi$ decay structure, we introduced a set of physics-motivated pairwise features — opposite-sign indicator, $\rho(770)$ mass window, vertex compatibility along the beam line, and charge product — injected as additive attention biases on top of the generic Lorentz features used by the Particle Transformer~\cite{particle_transformer_2022}. These features encode the $a_1 \to \rho \pi$ decay topology that generic kinematic features miss, and contribute a measurable improvement in end-to-end recall over an otherwise identical baseline.

We achieved state-of-the-art results for low-$p_T$ $\tau_h$ track finding, refining the initial pool of $>1100$ candidates down to a shortlist of approximately $50$ tracks while retaining the full pion triplet. The size of this shortlist is the relevant difficulty axis: the fewer tracks the cascade outputs, the more aggressively background must be suppressed, and the closer the output approaches the input requirements of downstream $\tau$-lepton reconstruction algorithms that operate on small fixed-size sets.

Finally, to support model development and reproducibility, we built a comprehensive evaluation suite — R@$K$ curves, P@$K$, and conditional recall stratified by $p_T$ and impact-parameter significance — and released the full training pipeline as open-source forks of the \texttt{weaver-core} and \texttt{particle\_transformer} frameworks, including network wrappers, training scripts, and HTCondor submission infrastructure for CERN lxplus.

\section{Structure Of The Thesis}

To start with, Chapter~2 reviews the theoretical background and prior work relevant to low-$p_T$ $\tau_h$ reconstruction — particle physics and the LHC trigger system, deep learning architectures for jet tagging (ParticleNet, Particle Transformer, ABCNet), and existing CMS approaches in the low-$p_T$ regime. Following it, Chapter~3 introduces the low-$p_T$ tau dataset used throughout the thesis — its feature set and per-track labelling scheme — together with the evaluation pipeline of recall-based and per-event metrics that anchors the comparisons in subsequent chapters. After that, Chapter~4 presents the proposed two-stage cascade reconstruction model, detailing the learned pre-filter, the Particle-Transformer-based reranker, and the physics-motivated pairwise attention biases. Subsequently, Chapter~5 reports the experimental evaluation — ablation studies, comparison against the existing CMS algorithm, and analysis stratified by $p_T$ and impact-parameter regimes — while Chapter~6 concludes the thesis and outlines limitations and directions for future work.